\section{Conclusions}
\label{sec:conclusions}

This article compared single-point evaluation of univariate polynomials in
three representations on an Apple M1 processor.  The Boolean-kernel basis
leads to a balanced reduction tree in which all folds at one level are
independent.  In the common cross-representation benchmark at
\(N=2^{20}\), this evaluator required \(1.0560\) ms, compared with
\(1.4116\) ms for the blocked parallel monomial implementation.  The
difference corresponds to \(25.2\%\) less wall time for the isolated
evaluation operation.

The arithmetic-backend experiment helps identify the source of this
difference.  Specialized Goldilocks multiplication and Montgomery
multiplication have almost identical isolated multiplication times.  The
fold operation does not: the specialized Goldilocks fold is faster than
its Montgomery counterpart, and the same ordering is observed in all five
independent runs of the complete kernel evaluator.  No similarly stable
ordering is observed for the blocked monomial evaluator.  The performance
difference is therefore associated with the combination of the
Boolean-kernel reduction structure and the arithmetic used inside the
fold, rather than with a universal advantage of one scalar multiplication
routine.

These measurements concern a single processor and an isolated polynomial
operation.  They do not establish an end-to-end speedup for a proof system.
Such a speedup depends on the representation used by the surrounding
protocol, the frequency of basis conversion, and the fraction of prover
time spent in the evaluated primitive.  A monomial-to-kernel conversion
requires \(O(N\log N)\) field additions, so repeated conversion can erase
the benefit measured here.

Further work should therefore address two questions.  The first is
architectural: the same evaluators should be measured on additional
out-of-order CPUs and, where appropriate, on SIMD and accelerator
implementations.  The second is algorithmic: the Boolean-kernel
representation should be integrated into a complete polynomial-commitment
or proof-system pipeline in which conversion, commitment, and opening costs
are measured together.  These experiments would determine whether the
single-point advantage observed here persists beyond the isolated
evaluation kernel.
